\RequirePackage{iftex}
\ifPDFTeX
\documentclass[a4paper,12pt]{article}
\usepackage[margin=24mm]{geometry}
\usepackage[T1]{fontenc}
\usepackage[utf8]{inputenc}
\usepackage{CJKutf8}
\else
\documentclass[a4paper,12pt]{jsarticle}
\usepackage[dvipdfmx,margin=24mm]{geometry}
\fi
\usepackage{amsmath,amssymb,amsthm}
\usepackage{enumitem}

\setlength{\parindent}{0pt}
\setlength{\parskip}{0.35em}
\setlist[itemize]{leftmargin=2em,itemsep=0.2em,topsep=0.3em}
\setlist[enumerate]{itemsep=0.4em,topsep=0.3em}

\newcommand{\feedbackqa}[2]{%
\item[]
\begin{tabular}[t]{@{}p{1.8em}@{\hspace{0.4em}}p{\dimexpr\linewidth-2.2em\relax}@{}}
Q. & #1 \\
A. & #2
\end{tabular}
\par\smallskip\hrulefill
}

\newcommand{\feedbacknote}[1]{%
\item[] #1\par\smallskip\hrulefill
}

\newcommand{\feedbackrule}{%
\item[] \leavevmode\hrulefill
}

\begin{document}
\ifPDFTeX
\begin{CJK}{UTF8}{min}
\fi

\begin{center}
{\LARGE 数学1A 第7回レジュメ}\\[0.4em]
偏微分の定義と計算
\end{center}


\section{二変数関数の復習}

二変数関数とは，2つの実数 $x,y$ に対して1つの実数 $f(x,y)$ を対応させる関数である。たとえば
\[
f(x,y)=x^2+y^2,\qquad
g(x,y)=xy+\sin y,\qquad
h(x,y)=e^{x+y}
\]
などが二変数関数である。


\section{偏微分係数}

関数 $f(x,y)$ の点 $(a,b)$ における $x$ に関する偏微分係数を
\[
\frac{\partial f}{\partial x}(a,b)
=\lim_{h\to 0}\frac{f(a+h,b)-f(a,b)}{h}
\]
で定義する。ただし，この極限が存在する場合に限る。

同様に，点 $(a,b)$ における $y$ に関する偏微分係数を
\[
\frac{\partial f}{\partial y}(a,b)
=\lim_{h\to 0}\frac{f(a,b+h)-f(a,b)}{h}
\]
で定義する。

$x$ に関する偏微分では $y$ を定数として固定する。$y$ に関する偏微分では $x$ を定数として固定する。

各点 $(x,y)$ で偏微分係数を考えると，次のような新しい関数が得られる。
\[
\frac{\partial f}{\partial x}(x,y)
=\lim_{h\to 0}\frac{f(x+h,y)-f(x,y)}{h},
\qquad
\frac{\partial f}{\partial y}(x,y)
=\lim_{h\to 0}\frac{f(x,y+h)-f(x,y)}{h}.
\]

これらを $f$ の偏導関数という。記号として
\[
f_x(x,y)=\frac{\partial f}{\partial x}(x,y),
\qquad
f_y(x,y)=\frac{\partial f}{\partial y}(x,y)
\]
とも書く。


\section*{4. 基本的な計算の考え方}

偏微分を計算するときは，一方の変数だけを変数と見て，他方を定数として扱う。

たとえば
\[
f(x,y)=x^3y^2+4xy-5y
\]
とする。$x$ に関して偏微分するとき，$y$ は定数であるから
\[
\frac{\partial f}{\partial x}(x,y)
=3x^2y^2+4y.
\]

$y$ に関して偏微分するとき，$x$ は定数であるから
\[
\frac{\partial f}{\partial y}(x,y)
=2x^3y+4x-5.
\]
\ifPDFTeX
\end{CJK}
\fi

\end{document}
